\documentclass[11pt,a4paper]{article}

% Packages
\usepackage[utf8]{inputenc}
\usepackage[T1]{fontenc}
\usepackage{amsmath, amssymb, amsthm}
\usepackage{graphicx}
\usepackage{booktabs}
\usepackage{hyperref}
\usepackage{geometry}
\usepackage{enumitem}
\usepackage{fancyhdr}
\usepackage{titlesec}
\usepackage{parskip}
\usepackage{xcolor}

% Page geometry
\geometry{margin=2.5cm, headheight=14pt}

% Color definitions
\definecolor{ImperialBlue}{RGB}{0, 62, 116}
\definecolor{DarkGray}{RGB}{51, 51, 51}
\definecolor{AccentOrange}{RGB}{214, 96, 77}

% Hyperref settings
\hypersetup{
    colorlinks=true,
    linkcolor=ImperialBlue,
    urlcolor=ImperialBlue,
    citecolor=ImperialBlue
}

% Header and footer
\pagestyle{fancy}
\fancyhf{}
\fancyhead[L]{\textcolor{ImperialBlue}{\small Blockchain Economics and Digital Assets}}
\fancyhead[R]{\textcolor{ImperialBlue}{\small Lecture 5 Notes}}
\fancyfoot[C]{\thepage}
\renewcommand{\headrulewidth}{0.4pt}

% Section formatting
\titleformat{\section}
  {\Large\bfseries\color{ImperialBlue}}
  {\thesection}{1em}{}
\titleformat{\subsection}
  {\large\bfseries\color{ImperialBlue}}
  {\thesubsection}{1em}{}
\titleformat{\subsubsection}
  {\normalsize\bfseries\color{DarkGray}}
  {\thesubsubsection}{1em}{}

% Custom commands
\newcommand{\keyword}[1]{\textbf{\textcolor{AccentOrange}{#1}}}

%----------------------------------------------------------------------------------------
\begin{document}

\begin{center}
    {\LARGE\bfseries\textcolor{ImperialBlue}{Blockchain Economics and Digital Assets}}\\[0.5em]
    {\Large Lecture 5: Stablecoins and Central Bank Digital Currencies}\\[1em]
    {\large Dr Daniele Bianchi}\\[0.5em]
    {\normalsize Queen Mary, University of London}\\[0.3em]
    {\normalsize Semester B, 2025/2026}
\end{center}

\vspace{1em}
\hrule
\vspace{1em}

\tableofcontents

\vspace{2cm}

%----------------------------------------------------------------------------------------
\section*{Overview}
\addcontentsline{toc}{section}{Overview}
%----------------------------------------------------------------------------------------

Bitcoin and other cryptocurrencies suffer from a fundamental problem as media of exchange: they are extraordinarily volatile. A currency that can move 10\% in a day is poorly suited for pricing goods, denominating loans, or settling contracts.

\keyword{Stablecoins} attempt to solve this problem by creating cryptocurrency tokens that maintain a stable value, typically pegged to the US dollar. They combine the programmability and accessibility of crypto with the price stability of traditional currencies. Stablecoins have become critical infrastructure: they exceed \$200 billion in market capitalisation, serve as the primary trading pairs on exchanges, and provide the stable unit of account that DeFi requires.

But stability mechanisms vary enormously in robustness. Some stablecoins are backed by real dollars; others by volatile cryptocurrency collateral; still others by algorithmic mechanisms with no backing at all. The collapse of Terra/Luna in May 2022---which destroyed \$60 billion in value within days---demonstrated that ``stable'' is not always accurate.

This lecture examines the economics of stablecoins and then turns to \keyword{Central Bank Digital Currencies} (CBDCs)---government-issued digital money that represents an alternative response to the same underlying demand for digital value transfer.

%----------------------------------------------------------------------------------------
\section{Why Stablecoins?}
%----------------------------------------------------------------------------------------

\subsection{The Volatility Problem}

Cryptocurrencies are volatile. Bitcoin routinely moves 5--10\% in a single day. This creates practical problems for payments (price risk between acceptance and conversion), lending (how do you price a loan in an asset that may halve?), and smart contracts (DeFi needs a stable unit of account).

\subsection{The Solution}

A stablecoin maintains a stable value relative to a reference asset, typically the US dollar. It retains cryptocurrency benefits---programmability, 24/7 availability, global accessibility---while eliminating volatility relative to the dollar.

The market exceeds \$200 billion, dominated by USDT (Tether, $\sim$\$150B) and USDC (Circle, $\sim$\$70B). Together they constitute roughly 90\% of stablecoin market cap.

%----------------------------------------------------------------------------------------
\section{Types of Stablecoins}
%----------------------------------------------------------------------------------------

\begin{center}
\begin{tabular}{llp{4cm}}
\toprule
\textbf{Type} & \textbf{Backing} & \textbf{Examples} \\
\midrule
Fiat-backed & USD in bank accounts, T-bills & USDT, USDC \\
Crypto-backed & Cryptocurrency (over-collateralised) & DAI \\
Algorithmic & No collateral; supply adjustments & UST (failed) \\
Commodity-backed & Gold, other commodities & PAXG \\
\bottomrule
\end{tabular}
\end{center}

The key distinction: off-chain collateral requires trusting the custodian; on-chain collateral is verifiable but volatile; algorithmic mechanisms have no fallback when confidence breaks.

%----------------------------------------------------------------------------------------
\section{Fiat-Backed Stablecoins}
%----------------------------------------------------------------------------------------

\subsection{The Basic Mechanism}

\textbf{Minting}: User sends \$1 to issuer $\rightarrow$ Issuer creates 1 stablecoin $\rightarrow$ Issuer deposits \$1 in reserves.

\textbf{Redemption}: User sends 1 stablecoin to issuer $\rightarrow$ Issuer burns the stablecoin $\rightarrow$ Issuer sends \$1 from reserves.

The peg is maintained by \keyword{arbitrage}. If market price exceeds \$1, arbitrageurs mint at \$1 and sell on market. If market price falls below \$1, they buy on market and redeem at \$1. This adjusts supply to maintain the peg.

\subsection{Tether (USDT) and USD Coin (USDC)}

Tether is the dominant stablecoin (\$150B+), launched in 2014 by Tether Limited (related to Bitfinex exchange).

\textbf{Controversies}:
\begin{itemize}
    \item No full audit---only ``attestations'' (point-in-time snapshots)
    \item Reserve quality improved over time: from 50\% commercial paper (2021) to 80\% T-bills (2024)
    \item 2021 CFTC settlement: \$41M for misrepresenting reserves
\end{itemize}

Despite controversies, USDT has maintained its peg through multiple crises. Users evidently trust redemption ability regardless of transparency concerns.

USDC is issued by Circle (in partnership with Coinbase), positioned as the transparent, regulated alternative:
\begin{itemize}
    \item Monthly attestations from Big Four accounting firms
    \item Reserves in regulated US institutions
    \item Circle is a licensed money transmitter
    \item Reserve composition: $\sim$80\% short-dated Treasuries, $\sim$20\% bank deposits
\end{itemize}

USDC was seen as ``safer''---until March 2023.

\subsection{Case Study: The USDC Depeg (March 2023)}

On March 10, 2023, Silicon Valley Bank collapsed. Circle disclosed that \$3.3 billion of USDC reserves (approximately 8\%) were held at SVB. Panic ensued: if Circle could not access these funds, redemptions at \$1 would be impossible.

USDC dropped to \keyword{\$0.87} on some exchanges over the weekend.

The resolution came when the US government guaranteed all SVB deposits on March 12. Circle confirmed full access to funds, and USDC returned to \$1 within 48 hours.

\textbf{Lessons}:
\begin{itemize}
    \item Even ``safe'' stablecoins face bank counterparty risk
    \item Transparency helped: Circle disclosed exposure quickly
    \item Government backstop restored confidence
    \item Fiat-backed stablecoins inherit traditional banking risks
\end{itemize}

\subsection{The Bank Run Problem}

Fiat-backed stablecoins face classic bank-run dynamics:

\begin{enumerate}
    \item A confidence shock triggers mass redemptions
    \item To meet redemptions, the issuer must liquidate reserves
    \item Fire sales may incur losses, reducing reserve coverage
    \item Reduced coverage increases fear, triggering more redemptions
    \item The feedback loop can break the peg even if reserves were initially adequate
\end{enumerate}

\textbf{Key difference from banks}: No deposit insurance, no central bank backstop. Mitigants include holding only the safest, most liquid assets (T-bills) and maintaining reserves $\geq$ 100\%.

%----------------------------------------------------------------------------------------
\section{Crypto-Backed Stablecoins}
%----------------------------------------------------------------------------------------

\subsection{The Mechanism}

If collateral is volatile cryptocurrency, how do you maintain a \$1 peg? The solution is \keyword{over-collateralisation}: require more collateral than stablecoins issued.

\textbf{Example}: To mint \$100 of stablecoins, deposit \$150 of ETH. If ETH drops 20\%, collateral is still worth \$120---sufficient to back the \$100 in stablecoins.

\textbf{Trade-offs}:
\begin{itemize}
    \item Transparent: Collateral is on-chain, verifiable by anyone
    \item No bank counterparty risk
    \item But capital inefficient: Need \$150 to create \$100
    \item Inherits crypto volatility risks
\end{itemize}

\subsection{DAI}

DAI is the leading crypto-backed stablecoin, launched in 2017 by MakerDAO. It is decentralised---no single company controls it---and DAI is essentially a loan against your crypto collateral.

\textbf{How it works}:
\begin{enumerate}
    \item Deposit 1 ETH (worth \$3,000) as collateral
    \item Minimum collateralisation ratio: 150\%
    \item Maximum DAI you can mint: \$3,000 / 1.5 = \$2,000
    \item Prudent users mint less (e.g., \$1,500 for 200\% collateralisation)
\end{enumerate}

\textbf{If ETH drops to \$2,000}:
\begin{itemize}
    \item Your collateral ratio falls to \$2,000 / \$1,500 = 133\%
    \item Below 150\% threshold $\rightarrow$ \keyword{liquidation risk}
    \item Liquidators can repay part of your debt and seize collateral at a discount
\end{itemize}

\textbf{Why this maintains the peg}:
\begin{itemize}
    \item DAI is always backed by $>$100\% collateral value
    \item Liquidations ensure undercollateralised positions are closed
    \item If DAI $<$ \$1: cheaper to repay debt, reducing supply
    \item If DAI $>$ \$1: profitable to mint more, increasing supply
\end{itemize}

\textbf{The USDC controversy}: DAI now accepts multiple collateral types, including USDC ($\sim$30\% of backing). This means DAI inherits USDC's risks. During the March 2023 USDC depeg, DAI also dropped briefly. There is tension between stability (accepting stable collateral) and decentralisation (avoiding centralised assets).

%----------------------------------------------------------------------------------------
\section{Algorithmic Stablecoins}
%----------------------------------------------------------------------------------------

\subsection{The Idea}

Maintain the peg through supply adjustments rather than collateral:
\begin{itemize}
    \item If price $>$ \$1: Increase supply (mint new coins) $\rightarrow$ price falls
    \item If price $<$ \$1: Decrease supply (remove coins) $\rightarrow$ price rises
\end{itemize}

The challenge: increasing supply is easy; decreasing supply is hard. Common mechanisms include seigniorage shares (issue ``bonds'' promising future stablecoins), dual-token systems (absorb volatility into a second token), or rebasing (automatically adjust balances).

All require \keyword{continued confidence}. If confidence breaks, there is no collateral to fall back on.

\subsection{Case Study: Terra/Luna (May 2022)}

\textbf{The mechanism}: UST (stablecoin) was paired with LUNA (volatile token). To mint 1 UST, burn \$1 worth of LUNA. To redeem 1 UST, burn UST and receive \$1 worth of LUNA. Arbitrage should maintain the peg.

\textbf{The sweetener}: Anchor Protocol paid \keyword{20\% APY} on UST deposits---unsustainably high yields that attracted massive capital.

\textbf{At peak} (April 2022): UST market cap $\sim$\$18 billion; LUNA market cap $\sim$\$40 billion; LUNA price $\sim$\$80.

\textbf{The death spiral} (May 7--13, 2022):
\begin{enumerate}
    \item Large UST withdrawals from Anchor triggered a confidence shock
    \item UST dropped slightly below \$1
    \item Holders redeemed UST for LUNA (minting new LUNA)
    \item LUNA supply increased $\rightarrow$ LUNA price fell
    \item Now \$1 of LUNA was worth less $\rightarrow$ needed \textit{more} LUNA per UST redeemed
    \item LUNA supply exploded: 350 million $\rightarrow$ 6.5 \textit{trillion}
    \item LUNA price: \$80 $\rightarrow$ \$0.0001
    \item With LUNA worthless, UST had no redemption value
    \item UST collapsed to \$0.10
\end{enumerate}

\textbf{Total losses}: Approximately \$60 billion destroyed in one week.

\subsection{Why Algorithmic Stablecoins Are Fragile}

The fundamental problem is \keyword{circular backing}:
\begin{itemize}
    \item UST's value depends on ability to redeem for LUNA
    \item LUNA's value depends on demand for UST
    \item If one loses confidence, both collapse
\end{itemize}

Contrast with fiat-backed stablecoins: USDC is backed by T-bills, which have value independent of USDC demand. External collateral breaks the reflexivity.

\textbf{Lesson}: Algorithmic stablecoins can work during good times. The question is whether they survive bad times. Terra did not. No algorithmic stablecoin has achieved both scale and long-term stability.

%----------------------------------------------------------------------------------------
\section{Stablecoin Risks and Regulation}
%----------------------------------------------------------------------------------------

\subsection{Summary of Risks}

\begin{center}
\begin{tabular}{lp{8cm}}
\toprule
\textbf{Risk} & \textbf{Description} \\
\midrule
Run risk & Mass redemptions exhaust reserves \\
Reserve quality & Illiquid or risky assets cannot meet redemptions \\
Counterparty risk & Bank failures, custodian problems \\
Operational risk & Smart contract bugs, key management failures \\
Concentration & Few stablecoins dominate; failure is systemic \\
Regulatory & Legal status uncertain; potential bans \\
\bottomrule
\end{tabular}
\end{center}

Stablecoins are deeply embedded in crypto markets. A major stablecoin failure could trigger cascading liquidations across DeFi.

\subsection{Regulatory Landscape}

\textbf{Europe (MiCA)}: The Markets in Crypto-Assets Regulation entered force in 2024. Stablecoin issuers must be authorised as electronic money institutions, hold 1:1 reserves in liquid assets, undergo regular audits, and have EU presence. Tether has been delisted from some EU exchanges; USDC is pursuing compliance.

\textbf{United States}: Fragmented, with no comprehensive stablecoin law. The SEC views some stablecoins as potential securities; the CFTC has brought enforcement actions; various Congressional bills propose bank-like regulation. Regulatory approach may shift significantly depending on political landscape.

%----------------------------------------------------------------------------------------
\section{Central Bank Digital Currencies}
%----------------------------------------------------------------------------------------

\subsection{What is a CBDC?}

\begin{quote}
A \keyword{Central Bank Digital Currency} is a digital form of central bank money, denominated in the national currency and issued as a direct liability of the central bank.
\end{quote}

A CBDC differs from:
\begin{itemize}
    \item \textbf{Physical cash}: Digital, not physical
    \item \textbf{Bank deposits}: Direct central bank liability, not commercial bank
    \item \textbf{Stablecoins}: Issued by government, not private companies
    \item \textbf{Bitcoin}: Centralised issuance, not decentralised
\end{itemize}

Currently, only banks have direct access to central bank money in digital form (reserves). A CBDC would give the public this access.

CBDCs can be divided into retail and wholesale CBDCs:

\textbf{Retail CBDC}: Available to households and businesses for everyday transactions. Competes with cash and bank deposits. Examples: China's Digital Yuan, Bahamas' Sand Dollar.

\textbf{Wholesale CBDC}: Available only to financial institutions for interbank settlements. Competes with existing reserve systems. Examples: Project Helvetia (Switzerland), various pilots.

Most policy discussion focuses on retail CBDCs, which raise the most significant questions.

\subsection{Why Are Central Banks Interested?}

\textbf{Advanced economies}:
\begin{itemize}
    \item Payment system resilience and efficiency
    \item Declining cash usage (especially in Nordic countries)
    \item Respond to private stablecoins
    \item Maintain monetary sovereignty
\end{itemize}

\textbf{Emerging economies}:
\begin{itemize}
    \item Financial inclusion (unbanked populations)
    \item Reduce cash handling costs
    \item Combat informal economy
    \item Improve cross-border payments
\end{itemize}

Different motivations lead to different designs.

\subsection{CBDC Design Choices}

\begin{center}
\begin{tabular}{lll}
\toprule
\textbf{Choice} & \textbf{Option A} & \textbf{Option B} \\
\midrule
Architecture & Direct (CB manages accounts) & Intermediated (banks distribute) \\
Technology & Account-based (identity) & Token-based (possession) \\
Interest & Interest-bearing & Non-interest-bearing \\
Privacy & Anonymous (like cash) & Traceable (like deposits) \\
Limits & Unlimited holdings & Capped holdings \\
\bottomrule
\end{tabular}
\end{center}

Most proposed CBDCs are intermediated (banks handle distribution) and non-anonymous (transactions traceable).

%----------------------------------------------------------------------------------------
\section{CBDC Trade-offs}
%----------------------------------------------------------------------------------------

\subsection{Bank Disintermediation}

\textbf{The concern}: If citizens can hold money directly at the central bank, why keep deposits at commercial banks?

\textbf{Potential consequences}:
\begin{itemize}
    \item Deposits flow from banks to CBDC
    \item Banks lose cheap funding source
    \item Credit availability may decrease
    \item During crises, ``digital bank runs'' could be instantaneous
\end{itemize}

\textbf{Proposed mitigants}:
\begin{itemize}
    \item Cap CBDC holdings (e.g., €3,000 per person---ECB proposal)
    \item Make CBDC non-interest-bearing (less attractive than deposits)
    \item Tiered remuneration (penalty rate above threshold)
\end{itemize}

\subsection{Privacy vs Surveillance}

CBDCs could offer cash-like privacy or enable unprecedented surveillance.

\textbf{Privacy concerns}:
\begin{itemize}
    \item Every transaction could be tracked by the state
    \item Spending patterns reveal sensitive information
    \item Potential for political targeting or social control
    \item ``Programmable money'' could restrict purchases
\end{itemize}

\textbf{Authority arguments}:
\begin{itemize}
    \item Traceability needed for anti-money laundering
    \item Tax compliance requires visibility
    \item Tiered privacy possible (anonymous for small amounts)
\end{itemize}

Design choices here reflect fundamental values about state-citizen relationships.

\subsection{Monetary Policy Implications}

CBDCs could enable new policy tools:
\begin{itemize}
    \item \textbf{Direct transmission}: Interest on CBDC goes straight to public
    \item \textbf{Negative rates}: Could charge for holding CBDC (harder with physical cash)
    \item \textbf{Helicopter money}: Direct transfers to citizens
    \item \textbf{Programmability}: Time-limited spending vouchers
\end{itemize}

Whether central banks should have these powers is a political question beyond technical feasibility.

\subsection{CBDCs vs Stablecoins}

\begin{center}
\begin{tabular}{lll}
\toprule
& \textbf{Stablecoins} & \textbf{CBDCs} \\
\midrule
Issuer & Private companies & Central bank \\
Backing & Reserves (varies) & Full faith of government \\
Regulation & Evolving & Government-controlled \\
Privacy & Pseudonymous & Design-dependent \\
Innovation & Rapid & Slow, deliberate \\
Risk & Issuer default & Political risk \\
\bottomrule
\end{tabular}
\end{center}

Whether CBDCs will replace stablecoins, coexist with them, or serve different use cases remains to be seen.

%----------------------------------------------------------------------------------------
\section{Summary and Looking Ahead}
%----------------------------------------------------------------------------------------

This lecture has covered stable value in the crypto ecosystem. Key takeaways:

\textbf{Stablecoins solve the volatility problem.} But stability mechanisms vary dramatically in robustness.

\textbf{Fiat-backed stablecoins dominate but require trust.} USDT is largest but controversial on transparency; USDC is more transparent but the March 2023 depeg showed bank counterparty risk.

\textbf{Crypto-backed stablecoins are transparent but capital-inefficient.} DAI requires over-collateralisation; liquidations maintain the peg.

\textbf{Algorithmic stablecoins are fragile.} Terra/Luna demonstrated catastrophic failure modes when confidence breaks.

\textbf{CBDCs offer government-backed digital money.} But raise significant trade-offs around disintermediation, privacy, and monetary policy power.

In the next lecture, we turn to digital ownership and tokenisation---representing assets as tokens on the blockchain, from speculative NFTs to institutional real-world asset tokenisation.

%----------------------------------------------------------------------------------------
\section*{Readings}
\addcontentsline{toc}{section}{Readings}
%----------------------------------------------------------------------------------------

\textbf{Required:}
\begin{itemize}
    \item Gorton, G. B., \& Zhang, J. (2023). ``Taming Wildcat Stablecoins.'' \textit{University of Chicago Law Review}, 90(3), 909--971. Economic analysis of stablecoin regulation.
\end{itemize}

\textbf{Supplementary:}
\begin{itemize}
    \item Bank for International Settlements. (2021). ``CBDCs: An Opportunity for the Monetary System.'' \textit{BIS Annual Economic Report}, Chapter III.
    \item Liu, Y., Tsyvinski, A., \& Wu, X. (2022). ``The Economics of Non-Fungible Tokens.'' Working paper. Section on stablecoin mechanisms.
    \item Uhlig, H. (2022). ``A Luna-tic Stablecoin Crash.'' NBER Working Paper 30256. Analysis of Terra/Luna collapse.
\end{itemize}

%----------------------------------------------------------------------------------------

\end{document}